% !TEX program = xelatex
% !TEX root = example.tex
% =====================================================================
%  EXAMPLE.TEX – ukázky syntaxe pro šablonu
%
%  Tento soubor je samostatný dokument. main.tex zůstává čistý a
%  připravený k psaní; sem se podívejte, když si potřebujete zkopírovat
%  obrázek, tabulku, rovnici, seznam, výpis kódu, citaci nebo odkaz.
%
%  Kompilace: Ctrl+Alt+B (LaTeX Workshop), task
%  "LaTeX: Build example (XeLaTeX 2x)" nebo přímo:
%  xelatex example.tex (dvakrát).
% =====================================================================
% !TEX root = ../main.tex
% =====================================================================
%  PREAMBULE – všechny balíčky, písmo a globální nastavení sazby.
%  Tento soubor se načítá z main.tex; běžně do něj nemusíte zasahovat.
% =====================================================================
\documentclass[12pt,a4paper]{article}

% ---------- geometrie stránky ----------
\usepackage[a4paper,
            left=2.5cm, right=2.4cm,
            top=2.0cm, bottom=2.5cm,
            footskip=16pt]{geometry}

% ---------- písmo a jazyk ----------
\usepackage{fontspec}
\usepackage{amsmath}
\usepackage{unicode-math}
\IfFontExistsTF{Times New Roman}
  {\setmainfont{Times New Roman}}
  {\setmainfont{TeX Gyre Termes}}
\IfFontExistsTF{Cambria Math}
  {\setmathfont{Cambria Math}}
  {\IfFontExistsTF{TeX Gyre Termes Math}
    {\setmathfont{TeX Gyre Termes Math}}
    {\setmathfont{Latin Modern Math}}}
\usepackage{polyglossia}
\setdefaultlanguage{czech}

% ---------- základní sazba ----------
\linespread{0.95}
\setlength{\parindent}{0.6cm}
\setlength{\parskip}{6pt}
\setlength{\emergencystretch}{3em}
\tolerance=1200
\hyphenpenalty=10000
\exhyphenpenalty=10000
\setlength{\abovedisplayskip}{8pt}
\setlength{\belowdisplayskip}{8pt}
\setlength{\abovedisplayshortskip}{4pt}
\setlength{\belowdisplayshortskip}{4pt}

\usepackage{xcolor}
\usepackage{graphicx}
\graphicspath{{obrazky/}{./}}
\usepackage{float}
\usepackage{needspace}
\usepackage[font=small,labelfont=bf,justification=centering,skip=6pt]{caption}
\usepackage{array}
\usepackage{multirow}
\usepackage{enumitem}
\usepackage{titlesec}
\usepackage{fancyhdr}
\usepackage{listings}
\usepackage[hidelinks]{hyperref}

% ---------- nadpisy ----------
\titleformat{\section}[hang]
  {\normalfont\fontsize{14}{16.8}\selectfont\bfseries}
  {\makebox[0.63cm][l]{\thesection.}}{0pt}{}
\titlespacing*{\section}{0pt}{14pt}{2pt}
\titleformat{\subsection}[hang]
  {\normalfont\fontsize{13}{15.6}\selectfont\bfseries}
  {}{0pt}{}
\titlespacing*{\subsection}{0pt}{12pt}{2pt}
\setcounter{secnumdepth}{1}

% ---------- seznamy ----------
\setlist[itemize]{leftmargin=1.27cm, labelindent=0.63cm,
  labelwidth=0.63cm, labelsep=0pt, align=left,
  topsep=2pt, itemsep=2pt, parsep=0pt, partopsep=0pt}
\setlist[enumerate]{leftmargin=1.27cm, labelindent=0.63cm,
  labelwidth=0.63cm, labelsep=0pt, align=left,
  topsep=2pt, itemsep=2pt, parsep=0pt, partopsep=0pt}

% ---------- zápatí ----------
\pagestyle{fancy}
\fancyhf{}
\fancyfoot[C]{- \thepage\ -}
\renewcommand{\headrulewidth}{0pt}
\renewcommand{\footrulewidth}{0pt}
\fancypagestyle{plain}{%
  \fancyhf{}%
  \fancyfoot[C]{- \thepage\ -}%
  \renewcommand{\headrulewidth}{0pt}%
  \renewcommand{\footrulewidth}{0pt}%
}

% ---------- české názvy ----------
\captionsetup[figure]{name=Obr.}
\captionsetup[table]{name=Tab.}
\renewcommand{\figurename}{Obr.}
\renewcommand{\tablename}{Tab.}
\renewcommand{\refname}{Použitá literatura}
\renewcommand{\contentsname}{Obsah}
\renewcommand{\listfigurename}{Seznam obrázků}
\renewcommand{\listtablename}{Seznam tabulek}
\setlength{\belowcaptionskip}{2pt}

% Některé názvy nastavuje polyglossia až na začátku dokumentu.
\AtBeginDocument{%
  \renewcommand{\refname}{Použitá literatura}%
  \renewcommand{\lstlistingname}{Výpis}%
}

% ---------- výpisy kódu ----------
\lstset{
  basicstyle=\ttfamily\small,
  breaklines=true,
  frame=single,
  rulecolor=\color{gray!50},
  backgroundcolor=\color{gray!5},
  keywordstyle=\bfseries,
  commentstyle=\itshape,
  stringstyle=\ttfamily,
  showstringspaces=false,
  tabsize=2,
  columns=fullflexible,
  keepspaces=true
}

% !TEX root = ../main.tex
% =====================================================================
%  VLASTNÍ MAKRA – zde jsou pomocné příkazy, které se používají v main.tex.
%  Běžně je nemusíte měnit.
% =====================================================================

% ---------- šedý text = místo k doplnění ----------
% Použití: \vstup{Text, který mám nahradit}
\newcommand{\vstup}[1]{\textcolor{gray!85!black}{\itshape #1}}

% ---------- obrázek s automatickým číslem a popiskem ----------
% Použití:
%   \obrazek[0.8\textwidth]{obrazek-1.png}{Popisek.}{obr:stitek}
%
% Argumenty:
%   1. (nepovinný) šířka obrázku, výchozí 0.85\textwidth
%   2. název souboru – hledá se v obrazky/ i v kořeni projektu
%   3. popisek pod obrázkem
%   4. štítek pro odkazy, např. \ref{obr:stitek}
%
% Dokud soubor neexistuje, vysází se šedý rámeček se jménem souboru.
\newcommand{\obrazek}[4][0.85\textwidth]{%
  \begin{figure}[H]
    \centering
    \setlength{\fboxsep}{0pt}%
    \IfFileExists{obrazky/#2}{%
      \includegraphics[width=#1,height=0.45\textheight,keepaspectratio]{obrazky/#2}%
    }{%
      \IfFileExists{#2}{%
        \includegraphics[width=#1,height=0.45\textheight,keepaspectratio]{#2}%
      }{%
        \fcolorbox{gray!45}{gray!8}{%
          \begin{minipage}[c][5cm][c]{\dimexpr#1-2\fboxrule\relax}
            \centering\color{gray}\itshape
            místo pro obrázek\\[4pt]
            \texttt{\detokenize{#2}}
          \end{minipage}}%
      }%
    }%
    \caption{#3}%
    \label{#4}%
  \end{figure}%
}


\begin{document}

\section*{Ukázky syntaxe}

\vstup{Tento soubor je samostatný dokument. Slouží jen jako referenční
přehled syntaxe – nic z~tohoto obsahu se nepřenáší do main.tex.
Příkazy si odsud zkopírujte do svého protokolu.}

\section{Nadpisy a text}

\subsection{Úrovně nadpisů}

Toto je běžný odstavec. Můžete v~něm použít \textbf{tučné písmo},
\textit{kurzívu} nebo \texttt{monospace}. V~českém textu se hodí
pevná mezera, například \verb|v~textu| nebo \verb|15.~11.~2026|.

\subsubsection{Jemnější nadpis}

Subsubsection je menší než subsection. V~šabloně se nečísluje, ale
může se hodit pro členění delší kapitoly.

\paragraph{Nadpis odstavce} Paragraph je ještě menší nadpis, který
zůstává na začátku odstavce.

Mezi speciální znaky patří \&, \%, \_, \#, \$ a~také české uvozovky
„takto“. Zpětné lomítko se v~textu píše jako \texttt{\textbackslash}.

\section{Obrázky}

Obrázek se vkládá příkazem \texttt{\textbackslash obrazek}. První
nepovinný argument je šířka, další tři jsou název souboru, popisek
a~štítek pro odkazy. Soubory patří do složky \texttt{obrazky/}.

% UKÁZKA – obrázek s vlastní šířkou
\obrazek[0.6\textwidth]{obrazek-1.png}{Ukázkový popisek obrázku.}{obr:ukazka}

Na obrázek se odkazujte jako na obr.~\ref{obr:ukazka}. Číslo i~popisek
se vytvářejí automaticky.

% UKÁZKA – druhý obrázek s výchozí šířkou 0.85\textwidth
% \obrazek{obrazek-2.png}{Popisek druhého obrázku.}{obr:druhy}

\section{Tabulky a rovnice}

\subsection{Tabulka}

Tabulku pište do prostředí \texttt{table} s~prostředím
\texttt{tabular}. Popisek patří nad tabulku.

% UKÁZKA – tabulka
\begin{table}[H]
  \centering
  \caption{Ukázková tabulka s~naměřenými hodnotami.}
  \label{tab:ukazka}
  \begin{tabular}{|l|c|r|}
    \hline
    \textbf{Parametr} & \textbf{Hodnota} & \textbf{Jednotka} \\ \hline
    Napětí & 5,00 & V \\ \hline
    Proud & 0,25 & A \\ \hline
  \end{tabular}
\end{table}

Odkaz na tabulku: tab.~\ref{tab:ukazka}.

\subsection{Rovnice}

Jedna číslovaná rovnice:

% UKÁZKA – rovnice
\begin{equation}
  y = a x + b
  \label{eq:ukazka}
\end{equation}

Odkaz na rovnici: \eqref{eq:ukazka}.

Více zarovnaných rovnic:

% UKÁZKA – align
\begin{align}
  y &= a x + b \label{eq:primka}\\
  z &= y^2     \label{eq:kvadrat}
\end{align}

Rovnice \eqref{eq:primka} a~\eqref{eq:kvadrat} jdou odkazovat zvlášť.

\section{Seznamy}

Odrážky:

% UKÁZKA – itemize
\begin{itemize}
  \item První odrážka.
  \item Druhá odrážka.
  \item Třetí odrážka.
\end{itemize}

Číslovaný seznam:

% UKÁZKA – enumerate
\begin{enumerate}
  \item První krok.
  \item Druhý krok.
  \item Třetí krok.
\end{enumerate}

Seznamy lze vnořovat.

\section{Kód}

Delší ukázku kódu vložte do prostředí \texttt{lstlisting}. Jazyk se
volí parametrem \texttt{language}, například \texttt{Python},
\texttt{C}, \texttt{C++}, \texttt{Java} nebo \texttt{SQL}.

% UKÁZKA – výpis kódu
\begin{lstlisting}[language=Python,
                   numbers=left,
                   caption={Ukázkový výpis kódu.},
                   label={lst:ukazka}]
def soucet(a, b):
    return a + b

vysledek = soucet(2, 3)
print(vysledek)
\end{lstlisting}

Odkaz na výpis: \ref{lst:ukazka}. Krátký název souboru pište přímo
jako \texttt{main.py}.

\section{Citace, poznámky a odkazy}

Literaturu citujte pomocí \texttt{\textbackslash cite\{...\}} a~seznamu
\texttt{thebibliography} na konci dokumentu.

% UKÁZKA – citace, poznámka pod čarou, URL a odkaz
Ukázková citace~\cite{ref:ukazka}. Tady je poznámka pod
čarou.\footnote{Toto je ukázková poznámka pod čarou.}
Webová adresa: \url{https://www.zcu.cz}. Odkaz s~vlastním textem:
\href{https://www.zcu.cz}{ZČU}.

Na stránku s~obrázkem se dostanete přes \texttt{\textbackslash pageref}:
str.~\pageref{obr:ukazka}.

\section{Zlomy a volitelné části}

Ruční zlom stránky:

\begin{verbatim}
\newpage
\end{verbatim}

Ochrana proti osamocenému nadpisu:

\begin{verbatim}
\needspace{8cm}
\section{Název kapitoly}
\end{verbatim}

Volitelný obsah, seznamy a přílohy:

\begin{verbatim}
\tableofcontents
\listoffigures
\listoftables
\appendix
\section{Přílohy}
\end{verbatim}

% =====================================================================
% POUŽITÁ LITERATURA
% =====================================================================
\renewcommand{\refname}{Použitá literatura}
\begin{thebibliography}{9}
  \bibitem{ref:ukazka}
  Autor. \textit{Název publikace}. Město: Vydavatel, 2024.
\end{thebibliography}

\end{document}
